\chapter{绪论}
\label{chap:introduction}
\raggedbottom

\section{研究背景及意义}
\label{sec:intro-background}

\subsection{医学影像智能分析及其临床价值}
\label{subsec:intro-medical-imaging-value}

医学影像将人体内部的解剖结构、组织形态及部分生理活动转换为可记录和比较的空间信息。X 射线、计算机断层成像、磁共振成像、超声成像和数字病理等技术具有不同的成像机理与信息尺度，但其临床解释通常都涉及对异常征象的识别、对器官或病灶范围的测量，以及对不同时间点、不同模态或不同个体影像的比较。医学影像智能分析因而不只是对图像给出一个类别标签，而是利用计算模型提取、组织和量化与诊疗任务相关的影像信息，为筛查、诊断、治疗规划和疗效随访提供可复核的分析结果。现有深度学习研究已覆盖分类、检测、分割、配准等多类医学影像任务\cite{litjens2017survey}。

从图像级分析看，医学图像分类主要建立图像内容与疾病类别、异常状态或风险水平之间的映射。根据任务粒度不同，模型可以处理整幅图像、局部候选区域或多幅影像组成的检查序列，输出离散类别或连续评分。这类结果可用于病例筛选、检查分诊和辅助鉴别，但其作用依赖于明确的适用人群、输入条件与决策阈值。以皮肤病变识别为例，Esteva 等利用临床图像和疾病标签训练卷积神经网络，说明端到端特征学习能够处理细粒度的医学图像分类问题\cite{esteva2017dermatologist}。此类研究展示了图像级预测的可行性，也表明分类模型的有效性与训练样本覆盖的疾病谱和成像条件密切相关。

从空间级分析看，医学图像分割需要为每个像素或体素赋予解剖结构、病灶或背景标签，其输出保留目标在图像空间中的位置、轮廓和形态。与图像级判断相比，分割结果能够支持器官体积、病灶负荷、形状变化和空间邻接关系的定量分析，并可作为手术规划、放射治疗勾画和纵向随访中的基础信息。U-Net 通过编码路径获取上下文信息，并利用对称解码路径恢复空间定位，形成医学图像分割中具有代表性的端到端建模方式\cite{ronneberger2015unet}。分割算法的实际价值并不只取决于平均重叠指标，还取决于轮廓误差、细小结构遗漏以及输出是否满足具体测量和规划要求。

进一步从图像间关系看，医学图像配准在两幅或多幅图像之间估计空间变换，使相同或对应的解剖位置建立一致关系。该任务可用于同一患者不同时间点影像的纵向比较、多模态信息融合、群体模板映射以及图像引导干预。传统配准方法往往针对每一对图像独立优化相似性度量与形变正则项；学习式方法则将图像对到形变场的映射参数化为神经网络，在训练完成后直接预测新的配准结果。VoxelMorph 是这一思路的代表性工作之一\cite{balakrishnan2019voxelmorph}。配准结果通常不直接给出疾病结论，但它为跨时间和跨模态分析提供统一坐标，是多阶段医学图像处理流程中的关键环节。

由此，分类、分割和配准分别形成图像级判断、空间级描述和图像间对应关系，三者在临床分析中具有互补性。例如，配准可以将不同检查对齐到可比较空间，分割可以量化局部结构及其变化，分类或风险预测再结合这些信息形成病例层面的判断。深度学习将许多依赖人工设计特征和逐例优化的流程转化为由数据学习的参数化映射，使模型能够在统一计算框架下处理高维医学图像。其作用主要体现在提供可重复的定量结果、减少部分重复性处理步骤，并扩大复杂影像信息被系统利用的范围。算法输出仍需结合临床资料、任务目标和人工复核进行解释，不能等同于独立的临床诊断。

不过，上述能力建立在训练数据能够充分描述目标任务的前提上。监督学习模型从图像及其标签中估计输入与输出之间的关系，样本规模、标注质量、疾病构成和采集条件会共同限定模型可以学习到的规律。医学影像数据的整理与专家标注本身具有较高成本，训练集也常难以兼顾规模、代表性和质量要求\cite{willemink2020preparing}。Zech 等对多医疗机构胸部 X 射线数据的研究进一步表明，模型在原训练机构内部的表现不能直接代表其在外部机构的表现，网络还可能利用与医院或采集流程相关的混杂信息\cite{zech2018variable}。因此，医学影像智能分析不仅要考虑固定数据集上的预测性能，还需要回答模型在真实数据访问条件下如何获得、保持和更新能力。

\subsection{医学影像数据、类别与分析目标的持续演化}
\label{subsec:intro-evolving-tasks}

医学影像模型面对的数据并非固定不变。患者人群、医疗中心、扫描设备、采集方案和重建方式会改变输入分布，疾病谱和医学语义类别会随新的诊断需求扩展，器官、模态、图像关系及分析目标也可能随研究阶段改变。传统离线学习通常在预先汇集的数据上完成一次优化，而持续学习要求模型在非平稳数据流中反复更新，并在吸收新知识的同时保持已经获得的能力\cite{hadsell2020embracing,vandeven2022three}。

本文关注三类变化：中心和采集条件改变引起的输入分布变化，医学图像类别的逐步增加，以及器官、模态和分析目标的改变。持续学习设定还取决于模型的输出方式和测试时是否提供任务身份，因此各章同时说明数据的到达方式与测试条件\cite{vandeven2022three,delange2022continual}。

任务类型说明模型预测什么，增量设定说明学习阶段之间哪些因素发生变化，训练条件则说明模型能够使用哪些信息。分割、分类和配准并不与三种增量设定一一对应，弱监督、联邦和回放条件也可以与不同设定组合。本文研究域变化与类别扩展下的弱监督持续分割、类别按批次到达的联邦持续分类，以及跨器官与跨模态的持续配准。

三类任务的输出和评价各有特点。弱监督分割在域增量中保持前列腺与背景的标签含义，在类别增量中逐步加入心脏结构；联邦分类采用测试任务身份已知的任务特定分类头；持续配准预测空间变换，任务变化来自器官、模态和配准关系。分割基准还包含跨器官任务增量序列，保留原研究中的 Organ-CL 标识，以便与其数据和图表对应。

顺序更新的核心困难之一是灾难性遗忘，即模型在学习后续阶段后，先前已经获得的能力下降\cite{delange2022continual}。遗忘与首次面对未见域时的泛化不足不同，也不能由训练结束后的单一平均性能充分描述。持续医学影像学习需要分别报告当前任务学习、历史任务保持、适用情况下的前向泛化，以及回放、参数增长、计算和通信等资源条件\cite{wang2024continualsurvey}。第二章将集中定义这些概念，第三章进一步把它们组织为 MedCL 平台设计和多维评价框架。

\subsection{医学影像数据的访问与标注约束}
\label{subsec:intro-incomplete-access}

常规集中式监督学习通常在已汇集的图像及其标签上训练，并允许重复读取训练样本。医学影像持续更新未必具备这些条件：新图像可能缺少与研究目标匹配的精细标签，旧任务数据可能无法在后续更新中重新访问，多中心原始影像也未必能够集中复制。因此，数据已经存在不等于其在当前训练阶段能够使用，模型实际可利用的信息往往少于潜在的数据资源\cite{willemink2020preparing,rieke2020future}。

造成这种差异的原因在于，医学数据的监督生成、时间访问和机构协作都受具体诊疗流程、研究任务与治理要求约束。影像采集不会自动产生与研究目标匹配的精细标签；历史样本是否可在后续优化中重访，取决于授权、存储和实验设定；跨机构原始影像能否汇集也不只是文件传输问题\cite{tajbakhsh2020imperfect,delange2022continual,rieke2020future,sheller2020federated}。这些约束分别影响监督信号、历史样本和跨中心信息的可用性，并可能同时出现。

本文将训练阶段无法完整调用监督信息、历史样本或跨中心原始数据的条件概括为训练信息受限。该概念描述信息的可用范围，不指图像中的随机缺失值，也不要求原始数据已被删除。下文分别说明三类约束如何改变持续更新时的学习条件。

\subsection{监督、历史访问与跨中心协作限制}
\label{subsec:intro-three-incomplete-dimensions}

部分监督直接改变模型能够获得的监督信号。分类任务可以使用检查级或病例级标签，但这类标签通常不能给出异常的空间位置和轮廓；分割任务需要逐像素或逐体素描述器官与病灶，配准研究也可能借助人工标志点或结构分割进行监督与评价。当标签由图像级细化到空间级时，复杂轮廓、小目标和多类别结构都会增加标注难度。图像级标签、点、框、涂鸦和部分切片标注能够降低监督获取负担，但它们不是密集掩膜的等价替代\cite{tajbakhsh2020imperfect}。

在涂鸦监督下，已标注像素提供局部类别位置，未标注区域的类别、轮廓和结构关系仍需由模型推断。若涂鸦覆盖集中在局部区域，或者不同类别获得的标注比例与完整区域差异较大，训练信号还会产生代表性偏差。部分交叉熵可以避免把未知像素直接视为背景，却不能单独恢复完整目标。有效的弱监督方法需要扩大可靠监督范围、校正类别与空间偏差，并利用与具体任务相符的结构信息。

历史数据的可访问程度决定后续更新能够使用哪些旧任务信息。联合训练可以混合所有已见任务的数据，有限回放只保留少量历史样本或其他表示，无原始样本回放方法则依赖旧模型、参数重要性或梯度子空间等摘要\cite{delange2022continual,wang2024continualsurvey}。比较方法时，需同时说明保存的对象、容量及其使用方式。

历史访问限制与任务演化需要区分。前者描述旧样本在当前优化中是否可用，后者描述当前输入分布、输出空间或任务目标是否发生变化。即使新旧数据来自相同分布，历史样本不可重访也会排除完整联合训练；反之，即使旧数据仍可访问，新的设备、模态、类别或器官任务也可能要求模型调整已有表征。明确这一区分，有助于解释有限回放与无回放方法的适用条件，也避免把“无回放”直接等同于“持续学习”。

多中心数据分散存放还限制了原始影像的集中使用。各中心能够在本地检查和使用影像，却不一定能够将其复制到统一服务器。联邦学习通过交换模型更新实现协同优化，因此需要处理客户端数据差异、参与变化、通信开销和服务器聚合等问题\cite{mcmahan2017fedavg,rieke2020future}。当各客户端还需要依次学习不同任务时，本地遗忘与跨客户端知识保持会相互影响，形成联邦持续学习问题。

上述三类约束改变模型可利用信息集合的不同部分。部分监督限制当前样本的标签内容；历史不可回放限制后续阶段重新使用旧样本的能力；分布式数据限制跨中心原始影像的集中访问。它们可以组合，但不能相互替代。例如，联邦训练不能补足涂鸦中缺失的目标轮廓，无回放设置也不能仅由多中心数据分散推出。方法和实验必须分别说明监督形式、历史访问和协作方式。

这些约束与任务演化共同决定了本文的研究问题：如何在当前监督有限、历史访问受限或数据分布于多个中心时持续学习，并使新获得的能力与既有能力得到分别评价。

常规集中式离线监督、元学习、持续学习和联邦学习分别从任务组织、时间到达、历史访问和原始数据位置等维度改变训练流程。它们不是互斥范式：元学习可以作为持续学习的优化工具，持续学习也可以在联邦环境中实施。图~\ref{fig:problem-settings-comparison}以常规集中式离线监督为参照，概括相关设定及其可组合关系。

\begin{figure}[htbp]
  \centering
  \captionsetup{font=small,skip=6pt}
  \begingroup
\linespread{1}\selectfont
\renewcommand{\arraystretch}{0.95}
\centering
\tikzset{
  >=Latex,
  frame/.style={draw=black, rounded corners=3pt, line width=0.65pt, fill=white},
  header/.style={draw=black, rounded corners=2pt, line width=0.65pt, fill=black!8,
    align=center, font=\thesisfigurefont\bfseries, minimum height=7mm},
  data/.style={draw=black, rounded corners=2pt, line width=0.55pt, fill=black!3,
    align=center, font=\thesisfigurefont, minimum height=7mm, inner xsep=4pt},
  process/.style={draw=black, rounded corners=2pt, line width=0.55pt, fill=white,
    align=center, font=\thesisfigurefont, minimum height=7mm, inner xsep=4pt},
  note/.style={align=center, font=\thesisfigurefont, text=black},
  flow/.style={-{Latex[length=1.7mm,width=1.1mm]}, draw=black, line width=0.55pt},
  exchange/.style={<->, draw=black, line width=0.55pt}
}

\resizebox{\thesisfigurewidth}{!}{%
\begin{tikzpicture}[x=1.65cm,y=0.74cm]

\begin{scope}[shift={(0,16.35)}]
  \draw[frame] (0,0) rectangle (8.35,4.75);
  \node[header, minimum width=7.75cm] at (4.175,4.28)
    {参照：常规集中式离线监督};
  \node[data, minimum width=2.05cm] (s-train) at (1.45,2.92)
    {预先汇集的\\训练数据};
  \node[process, minimum width=1.70cm] (s-opt) at (4.05,2.92)
    {一次性优化};
  \node[process, minimum width=1.65cm] (s-model) at (6.75,2.92)
    {固定模型};
  \node[data, minimum width=2.05cm] (s-test) at (1.45,1.13)
    {独立测试样本\\相对稳定条件};
  \node[process, minimum width=1.70cm] (s-pred) at (6.75,1.13)
    {预测与评价};
  \draw[flow] (s-train) -- (s-opt);
  \draw[flow] (s-opt) -- (s-model);
  \draw[flow] (s-test) -- (4.05,1.13) -- (s-pred);
  \draw[flow] (s-model) -- (s-pred);
\end{scope}

\begin{scope}[shift={(0,10.90)}]
  \draw[frame] (0,0) rectangle (8.35,4.75);
  \node[header, minimum width=7.75cm] at (4.175,4.28)
    {元学习：任务组织与快速适应};
  \node[data, minimum width=1.95cm] (m-source) at (1.30,2.93)
    {多个源任务\\$\mathcal{T}_1,\ldots,\mathcal{T}_K$};
  \node[process, minimum width=1.65cm] (m-meta) at (4.05,2.93)
    {元训练};
  \node[process, minimum width=1.80cm] (m-init) at (6.80,2.93)
    {可适应初始化};
  \node[data, minimum width=1.95cm] (m-support) at (1.30,1.13)
    {新任务少量数据\\支持集};
  \node[process, minimum width=1.65cm] (m-adapt) at (4.05,1.13)
    {任务内适应};
  \node[process, minimum width=1.80cm] (m-task) at (6.80,1.13)
    {新任务模型};
  \draw[flow] (m-source) -- (m-meta);
  \draw[flow] (m-meta) -- (m-init);
  \draw[flow] (m-support) -- (m-adapt);
  \draw[flow] (m-init.south) -- ++(0,-.4) -| (m-adapt.north);
  \draw[flow] (m-adapt) -- (m-task);
\end{scope}

\begin{scope}[shift={(0,5.45)}]
  \draw[frame] (0,0) rectangle (8.35,4.75);
  \node[header, minimum width=7.75cm] at (4.175,4.28)
    {持续学习：时间到达与历史访问};
  \node[data, minimum width=1.15cm] (c-t1) at (0.95,2.85) {$\mathcal{T}_1$};
  \node[process, minimum width=1.15cm] (c-m1) at (2.45,2.85) {$\theta_1$};
  \node[data, minimum width=1.15cm] (c-t2) at (3.95,2.85) {$\mathcal{T}_2$};
  \node[process, minimum width=1.15cm] (c-m2) at (5.45,2.85) {$\theta_2$};
  \node[process, minimum width=1.15cm] (c-mt) at (7.60,2.85) {$\theta_t$};
  \draw[flow] (c-t1.east) -- (c-m1.west);
  \draw[flow] (c-m1.east) -- (c-t2.west);
  \draw[flow] (c-t2.east) -- (c-m2.west);
  \draw[flow] (c-m2.east) -- (6.20,2.85);
  \node[note, inner xsep=0pt] at (6.48,2.85) {$\cdots$};
  % 箭头在省略号之后、theta_t 方框之前结束，避免箭头侵入方框。
  \draw[flow] (6.72,2.85) -- ([xshift=-0.8mm]c-mt.west);
  \node[data, minimum width=2.55cm] (c-history) at (2.05,1.35)
    {可访问的历史信息\\完整／有限／无回放};
  \node[process, minimum width=2.55cm] (c-eval) at (6.20,1.35)
    {持续评价\\旧任务、当前任务与未来任务};
  \draw[flow] (c-history) -- (4.15,1.35) -- (c-eval);
  \draw[flow] (c-mt.south) -- ++(0,-.35) -| (c-eval.north);

\end{scope}

\begin{scope}[shift={(0,0)}]
  \draw[frame] (0,0) rectangle (8.35,4.75);
  \node[header, minimum width=7.75cm] at (4.175,4.28)
    {联邦学习：数据位置与协同优化};
  \node[data, minimum width=3.7cm] (f-c1) at (1.55,3.10)
    {客户端 1：本地数据};
  \node[data, minimum width=3.7cm] (f-c2) at (1.55,2.03)
    {客户端 2：本地数据};
  \node[note] at (1.55,1.38) {$\vdots$};
  \node[data, minimum width=3.7cm] (f-ck) at (1.55,0.70)
    {客户端 $K$：本地数据};
  \node[process, minimum width=2.05cm, minimum height=2.2cm] (f-server) at (4.65,2.03)
    {服务器\\聚合与下发};
  \node[process, minimum width=1.95cm] (f-global) at (7.10,2.03)
    {共享模型};
  % 各客户端连接独立端口，三条双向线等长且水平平行。
  \draw[exchange] (f-c1.east) -- (f-server.west |- f-c1.east);
  \draw[exchange] (f-c2.east) -- (f-server.west);
  \draw[exchange] (f-ck.east) -- (f-server.west |- f-ck.east);
  \draw[exchange] (f-server) -- (f-global);

\end{scope}


\end{tikzpicture}%
}
\endgroup

  \caption[与本文相关的学习设定及其关系。]
  {与本文相关的学习设定及其关系。常规集中式离线监督在预先汇集的训练数据上完成一次性优化，并在相对稳定的测试条件下评价；元学习主要改变任务组织和快速适应方式；持续学习强调数据按时间到达、模型连续更新及历史信息的可用条件；联邦学习强调原始数据分布于多个客户端并通过双向模型通信协同优化。三种设定描述不同维度，可以相互组合。}
  \label{fig:problem-settings-comparison}
\end{figure}

\section{医学影像持续学习研究现状}
\label{sec:intro-related-work}

\subsection{面向数据分布演化的医学影像持续学习}
\label{subsec:intro-medical-continual-learning}

持续学习研究关注模型在数据或任务顺序到达时如何累积知识，并在给定历史访问和资源条件下维持既有能力。相关综述通常从回放、正则化、参数隔离及优化或表示约束等路线组织方法，并以稳定性--可塑性及任务内、任务间泛化刻画其目标\cite{delange2022continual,wang2024continualsurvey}。医学影像中的顺序变化还可能涉及样本批次、类别、器官、医疗中心、设备、采集参数和模态；任务身份、标签覆盖、输出结构与测试方式会进一步改变问题定义\cite{kumari2025medicalcl}。因此，研究现状需要一并核对任务序列、历史访问和评价方式，而不能只按方法名称比较。

通用方法保留的历史信息及其资源代价并不相同。iCaRL、弹性权重巩固和 Learning without Forgetting 分别体现样本回放、参数正则和函数正则的代表性思想\cite{rebuffi2017icarl,kirkpatrick2017ewc,li2018lwf}。这些机制在医学影像中的适用性还受到输入维度、输出结构和历史访问条件的影响。原始方法主要建立在分类设定上，其方法分类不能直接说明在高分辨率体数据、密集预测或特定访问条件下的效果。

早期持续医学图像分割研究主要关注设备和采集参数变化。Karani 等将不同扫描仪和采集参数视为顺序到达的相关域，共享卷积特征并保留域特定批归一化参数\cite{karani2018lifelong}。该设计减少新域适配对旧域的直接干扰，但需要为域保留特定状态，并依赖域身份或相应参数选择。此后，S3R 根据分割形状和语义估计参数重要性，对跨中心持续分割中的选择性参数变化进行约束\cite{zhang2023s3r}。这类方法把密集预测输出的结构属性纳入知识保持，但保护强度仍需在新中心适应和旧中心保持之间调节。

进一步的标准化框架研究表明，分类领域的结论不能直接套用于医学分割。Lifelong nnU-Net 在统一分割管线中比较回放、EWC、LwF、Riemannian Walk 和 MiB，并在多个持续分割用例中跟踪旧任务和新任务性能\cite{gonzalez2023lifelong}。该研究在其测试条件下观察到，所比较方法没有持续获得正向后向迁移，回放的遗忘相对较少，而任务特定输出头的影响有限。这些结论属于指定网络、数据序列和方法集合，不能扩展为所有医学分割场景的一般排序。

持续语义分割还具有分类任务中不完全相同的监督变化。MiB 指出，当当前阶段只标注新增类别时，旧类别像素可能被并入背景，导致背景语义随阶段改变，并据此调整损失与蒸馏\cite{cermelli2020mib}。PLOP 进一步利用多尺度特征蒸馏和背景伪标签维持旧类别信息\cite{douillard2021plop}。相关综述将这一现象概括为持续语义分割中的背景语义漂移，并按数据回放和无旧数据路线整理已有方法\cite{yuan2024continualseg}。对于医学分割，该问题是否出现取决于标签设定；在所有结构都被完整标注的域增量场景中，不能机械套用类别增量背景漂移结论。

当分割目标随阶段扩展时，研究开始使用器官特定组件和伪标签保护既有结构。Zhang 等在腹部多器官与肿瘤分割中采用器官特定输出头，并利用旧模型产生的伪标签辅助顺序学习\cite{zhang2023continualorgan}。这类设定可表现为新增解剖类别，也可因任务身份和输出组织方式被视为新任务。因此，“新增器官”本身不足以唯一决定增量类型，必须说明共享输出空间、测试时上下文和累计预测要求。

现有医学持续分割研究采用的场景和评价方式仍较分散。跨扫描仪或采集参数的顺序学习通常保持标签空间不变，类别扩展或跨器官任务研究则可能改变累计输出空间；不同工作对任务身份、输出头、旧类标签处理、历史数据访问和模型容量的设定也不一致\cite{karani2018lifelong,zhang2023s3r,gonzalez2023lifelong,zhang2023continualorgan}。在这些条件没有被一并说明时，最终平均性能、旧任务变化与当前任务学习结果难以直接横向比较。

总体来看，医学影像持续学习已从通用抗遗忘机制的直接迁移，发展到面向跨中心域变化、类别或器官扩展和密集预测语义变化的专门设计。现有方法分别利用回放、参数或函数正则、任务特定结构、伪标签和特征约束，但其结论受任务序列、历史访问、测试上下文和评价维度影响。如何在医学含义明确的设定下一并考察知识保持、当前任务学习、前向泛化与资源代价，仍是该方向需要解决的共性问题\cite{kumari2025medicalcl,yuan2024continualseg}。

本文以跨中心分割、心脏结构类别扩展和跨器官分割构建三类基准。后续章节在说明测试任务身份和输出方式的基础上，分别研究弱监督分割、联邦分类和持续配准。

\subsection{面向标签与监督演化的弱监督持续分割}
\label{subsec:intro-weak-partial-supervision}

弱监督医学图像分析旨在减少对完整像素级或体素级标注的依赖，并保留完成目标任务所需的监督信息。按照标注所提供的空间范围，已有研究使用图像级标签、点、包围框、涂鸦和部分切片标注等不同形式\cite{tajbakhsh2020imperfect}。这些标注的获取负担和信息粒度并不相同：图像级标签主要说明目标是否存在，点和涂鸦提供局部类别位置，包围框限定大致区域，而密集掩膜进一步描述完整轮廓和内部结构。因此，弱监督医学图像分割的核心不是简单减少标签数量，而是从不完整监督中恢复未被直接标出的空间信息。

涂鸦监督早期在自然图像语义分割中形成了标签传播的基本路线。ScribbleSup 将涂鸦标记、图像外观和网络预测纳入图模型，传播标签至未标注像素并更新分割网络\cite{lin2016scribblesup}。面向医学图像，Can 等比较不同训练策略，并在心脏和前列腺分割中直接利用涂鸦像素提供监督\cite{can2018scribble}。选择性损失或部分交叉熵只在已标注位置计算分类误差，能够避免把未知像素错误地当作背景，但未标注区域只能通过共享表征和网络归纳偏置受到间接影响。

一类研究通过正则项直接约束未标注区域，而不先生成完整硬标签。Tang 等将成对关系和条件随机场等结构约束写入弱监督分割损失，使图像信息参与未标注像素的优化\cite{tang2018regularized}。Kervadec 等进一步将目标区域大小等全局不等式约束转化为可微惩罚项，用领域知识限定网络输出\cite{kervadec2019constrained}。这类方法避免把单次预测直接固定为伪真值，但其有效性取决于约束是否与具体器官、病灶和数据分布相符；过强或不准确的先验同样可能限制模型学习。

另一类研究通过多分支预测或教师模型构造伪标签，以增加被监督的像素。Luo 等采用双分支网络，并动态混合两个分支的预测作为辅助监督\cite{luo2022scribble}。DMSPS 使用动态混合的软伪标签缓解硬伪标签的过度置信问题，并依据低不确定性预测扩展稀疏涂鸦\cite{han2024dmsps}。伪标签能够扩大监督覆盖，但其信息来自当前模型或相关分支；若早期预测存在系统性偏差，错误仍可能在后续训练中被强化。

针对体数据和结构信息，已有方法还从切片关联、轮廓和特征表示等角度补充监督。Scribble2D5 将相邻切片信息、标签传播和轮廓预测结合，用于涂鸦监督的体数据分割\cite{chen2022scribble2d5}。Zhou 等利用超像素引导涂鸦传播，并通过类别级对比正则增强类内一致性和类间区分\cite{zhou2023scribblewalking}。这些方法说明，稀疏标签之外的图像结构和表示关系可以提供附加约束；但超像素质量、切片连续性和特征紧致性并非在所有医学目标上都同样可靠。

总体而言，已有涂鸦监督方法主要沿标签传播、结构正则、伪标签、多分支一致性、跨切片关系和对比表示等路线扩大有限标注的有效作用范围\cite{lin2016scribblesup,tang2018regularized,kervadec2019constrained,luo2022scribble,han2024dmsps,chen2022scribble2d5,zhou2023scribblewalking}。这些方法在监督来源、先验假设和噪声控制方式上不同，其有效性也依赖目标结构、涂鸦质量和训练数据分布。

然而，现有弱监督分割研究多数在固定数据分布和静态标签空间中训练与评价，而持续分割研究通常假设当前阶段具有较完整的分割标签。任务顺序、稀疏监督、旧类标签处理与历史访问叠加出现时，伪标签既可能扩展当前监督，也可能携带由旧模型偏差或遗忘产生的错误。因而，部分监督与持续学习的叠加不能由静态弱监督结果或密集监督持续学习结果直接推出，仍需在明确任务设定和标注预算下单独研究。


当新中心图像持续到达，而训练标注只有稀疏涂鸦时，监督不足、域变化和遗忘会同时影响分割结果。第四章在密集监督分割基准的六中心前列腺和三任务心脏结构序列上进一步研究涂鸦监督，分别考察中心适应、旧域保持和类别扩展下的分割表现。

\subsection{面向类别扩展与多中心协作的联邦持续学习}
\label{subsec:intro-federated-continual-learning}

联邦学习通过在各参与方本地训练并在服务器端聚合更新，使多个医疗中心能够在不集中原始影像的条件下协同建模\cite{mcmahan2017fedavg,rieke2020future,sheller2020federated}。既有医学联邦学习主要研究固定任务下的数据异质性、优化收敛和跨中心泛化；当各中心还需要依次学习新的类别、域或诊断任务时，客户端本地数据流和服务器聚合过程都具有时间维度，问题由静态联邦学习扩展为联邦持续学习。

联邦持续学习要求各客户端从私有任务序列中更新模型，并通过通信形成能够覆盖多个客户端历史知识的共享模型。Yang 等将现有系统概括为同步和异步两类：同步设定中客户端按照共同任务顺序推进，异步设定中各客户端的任务内容、顺序或进度可以不同\cite{yang2024fclsurvey}。这一划分有助于说明服务器何时接收知识，但不能替代对标签空间、客户端参与、任务身份和历史数据访问条件的具体说明。

该问题包含两类相互耦合的知识损失。客户端在学习后续任务时可能降低旧任务性能，形成时间维度的灾难性遗忘；服务器聚合异质客户端更新时，又可能覆盖只存在于部分中心的任务知识。Yang 等将后一现象类比为空间维度的遗忘\cite{yang2024fclsurvey}。相关表述强调联邦聚合对跨客户端历史知识的影响，其具体评价仍取决于任务序列、共享模型和测试对象的定义。

早期联邦持续学习方法主要从参数分解和跨客户端知识选择入手。FedWeIT 将模型分为全局共享参数和稀疏任务特定参数，并根据任务相关性选择性接收其他客户端知识\cite{yoon2021fedweit}。这种设计可以减少不相关任务之间的直接干扰，但需要保存和管理任务相关状态，其效果也依赖任务关系估计和客户端通信条件。

与参数分解不同，知识蒸馏提供了另一类融合方式。CFeD 在客户端和服务器两侧使用无标签代理数据进行蒸馏，并让部分客户端学习新任务、部分客户端复习旧任务\cite{ma2022cfed}。在类别增量场景中，GLFC 结合类别感知梯度补偿、类别关系蒸馏、代理服务器和原型梯度通信，分别缓解本地类别不平衡与全局模型遗忘\cite{dong2022glfc}。这类方法减少了对完整旧图像的直接访问，但引入代理数据、旧模型选择或原型构造等附加假设。

Ke 等提出 TPPR，将视觉提示、任务感知梯度投影和提示回放结合，用于参数高效的联邦持续学习\cite{ke2025tppr}。该方法依托预训练模型，将更新集中于少量可训练参数；本文第五章则使用从头训练的医学分类网络，研究梯度子空间的跨客户端融合。

从更广的知识融合角度看，已有研究把相关方法整理为回放、聚类、参数正则、参数隔离、动态结构、原型和知识蒸馏等路线\cite{yang2024fclsurvey}。不同路线传递的数据、模型或输出信息不同：回放方法保留样本或替代数据，参数方法直接约束或选择模型更新，原型和蒸馏方法则以压缩表示或预测关系传递知识。

同时，联邦环境中的信息交换还需要与隐私机制区分。原始数据留在本地并不意味着共享梯度、参数、特征或原型不含训练信息，使用公共代理数据也会改变适用条件\cite{rieke2020future,yang2024fclsurvey}。因此，评价联邦持续医学影像方法时，应报告传输对象、历史数据或替代信息、可信参与方和附加保护机制，而不能仅以“数据不出中心”宣称形式化隐私。

医学影像中的联邦持续研究开始共同处理空间异质性和时间变化。DynBC 在病理图像分割中使用公共参考数据评价候选更新的连续性，以联合缓解客户端漂移和顺序学习遗忘\cite{babendererde2025dynbc}。该方法表明，服务器可以利用独立参考信息筛选跨中心更新；其适用性取决于公共参考数据能否获得，以及这些数据能否代表目标分布变化。

在医学图像分类中，Zhou 和 Wang 将多种持续学习算法与联邦训练结合，用于顺序胸部感染类别识别，并比较不同网络结构在其公开数据设定下的表现\cite{zhou2025thoracicfcl}。其客户端和中心由公开数据模拟，相关结果反映该实验设置下的表现。近期研究还将显式隐私机制加入联邦持续医疗分类。DP-FedEPC 在胸部 X 射线分类中结合弹性权重巩固、潜在原型和客户端差分隐私随机梯度下降，并报告对应隐私预算\cite{sinhal2026dpfedepc}。其中，形式化差分隐私来自梯度裁剪、噪声机制和隐私核算，而不是联邦聚合本身。

无回放方法试图在不保存旧图像或生成样本的条件下保护历史任务。梯度投影记忆利用旧任务的表示构建重要子空间，并将新任务梯度投影到该子空间的正交补\cite{saha2021gpm}。该方法最初面向观察完整任务序列的单一学习器；在联邦环境中，各客户端独立构建的保护子空间只能覆盖本地历史，无法直接包含仅由其他中心学习的任务知识。如何在服务器端汇集跨客户端历史保护信息，并避免无差别保护对新任务造成过强约束，是无回放联邦持续医学影像学习需要进一步解决的问题。

总体来看，联邦持续学习已从参数分解、蒸馏和类别增量补偿，扩展至更新筛选、显式差分隐私和无回放梯度空间保护。现有研究的任务顺序、同步方式、公共数据依赖、回放形式、通信对象和隐私假设差异较大。对于多中心无回放医学图像分类，仍需回答跨客户端历史知识如何形成全局保护，以及统一保护如何避免压低异质客户端的新任务学习能力。

第五章研究类别批次持续到达、客户端分布不一致且历史原始样本不回放时的联邦分类。测试时提供任务身份，并使用对应的任务特定分类头；训练中保留模型与低秩梯度子空间摘要，用于保护跨客户端历史知识。

\subsection{面向分析目标演化的持续医学图像配准}
\label{subsec:intro-registration-cl}

医学图像配准从固定图像和移动图像中估计连续空间变换，其输出与分类、分割中的离散语义标签不同。传统学习式配准通常针对相对固定的器官、模态和图像关系构建训练集，并在训练完成后直接预测新的形变场\cite{balakrishnan2019voxelmorph}。当脑部 MR、腹部 CT、肺部 CT 及腹部 MR--CT 等配准任务顺序到达时，输入外观、解剖区域、模态组合和形变分布均可能改变，单一网络的后续更新因而可能降低此前任务上的空间匹配能力\cite{wang2024samcl}。

持续配准不能机械套用分类任务的类别空间。其任务变化来自医学分析目标及配准关系的改变，历史能力则需要通过 Dice、目标配准误差及形变相关指标分别评价。有限经验回放可以保留部分历史图像对，但三维图像及形变模式的复杂性使少量样本难以覆盖旧分布；过强的历史约束又会压缩新任务的更新空间。元学习为跨任务快速适应提供初始化，锐度感知优化进一步关注参数邻域中的损失变化，但二者在持续环境中的作用仍受任务顺序和回放条件影响\cite{hadsell2020embracing,wang2024samcl}。

持续配准需要同时考察旧任务精度、新任务适应及模型对未参与训练数据的直接表现。第六章在四项配准任务和有限经验回放条件下研究这些问题。

\section{关键科学与工程问题}
\label{sec:intro-challenges}

分割、分类和配准虽然具有不同的预测对象和训练条件，但都需要判断模型在顺序更新中学到了什么、遗忘了什么。围绕评测平台、能力评价及三类方法，本文重点讨论以下五个问题。

\subsection{多任务研究的组织与评测流程}
\label{subsec:intro-challenge-platform}

不同任务的模型输出和测试参照各不相同，怎样将阶段模型、任务顺序、数据版本和评价结果准确对应，是组织实验时首先需要解决的问题。第三章据此设计 MedCL 平台，通过模型或预测提交、配置检查、任务评分和结果展示管理评测过程，同时保留各任务自己的模型结构与指标。

\subsection{持续学习能力的多维评价}
\label{subsec:intro-challenge-evaluation}

训练结束后的单一平均性能无法区分当前任务没有学好、旧任务被遗忘，或模型对未来任务缺乏直接泛化。分割可采用 Dice、IoU 和边界指标，分类采用准确率、AUC 或 F1，配准还可能同时使用 Dice、毫米级 TRE 和形变质量指标；这些量纲和优劣方向不同，不能直接平均为统一综合分数。评价框架需要在保留任务性能指标的同时，记录当前学习、历史保持、遗忘或后向迁移、适用情况下的前向泛化，以及存储、计算和通信条件。第三章给出共同的阶段--任务矩阵与指标定义，第四章通过持续分割基准开展相应的实证分析。

\subsection{稀疏监督下的跨中心分割与历史保持}
\label{subsec:intro-challenge-scribble}

新的中心和扫描条件会改变图像分布，而涂鸦只覆盖少量像素，不能直接给出完整区域与轮廓。当前阶段需要从稀疏监督中恢复结构，后续阶段还需防止模型丢失旧中心能力。扩大监督范围可能引入伪标签噪声，强化历史约束又可能妨碍对新域的适应。第四章由此在域增量与类别增量设置中研究弱监督和有限回放：以部分交叉熵、全局一致性和空间先验增强当前监督，并以历史图像--涂鸦回放、参考特征约束和回放弱监督优化保护旧域。

\subsection{无回放联邦分类中的跨客户端知识保持}
\label{subsec:intro-challenge-fedsubmerge}

多中心医学分类中，新的疾病或医学语义类别可能持续出现，原始影像又不能默认集中。客户端仅依据当前本地数据更新时，本地类别比例、样本量和图像来源差异会使更新偏离联邦总体历史；无历史原始样本回放进一步排除了集中重训。单个客户端构建的保护信息只能覆盖自身经历，无差别汇总全部历史方向又可能过度限制当前学习。第五章因此研究如何以低秩主梯度子空间汇集跨客户端历史，并通过统一或客户端相关的保护空间协调全局稳定性与本地可塑性。

\subsection{任务变化下配准模型的适应与保持}
\label{subsec:intro-challenge-samcl}

配准任务随器官、模态和配准关系改变时，模型输出仍是连续形变场，而不是逐步增加的离散类别。有限历史图像对只能近似旧任务分布，顺序更新既可能损害历史配准能力，也可能因过强约束无法学习新的形变模式。第六章研究如何将有限经验回放、一阶元持续学习和锐度感知更新结合，在固定任务序列中协调历史保持、当前适应以及预设域内与跨任务测试中的直接表现。


\section{本文主要研究内容}
\label{sec:intro-research-content}

本文先建立评测平台，再在具体医学影像任务中研究训练信息受限时的持续学习方法。第三章给出 MedCL 的任务组织和评价设计；第四章构建持续分割基准，并在相同病例划分与测试设置下研究弱监督分割；第五、六章分别研究无回放联邦分类和有限回放配准。三类任务使用不同的性能指标，但都通过阶段矩阵分析历史保持和当前学习，并记录模型扩展及历史信息使用。上述五个问题记为 RQ1 至 RQ5，与研究内容的对应关系见表~\ref{tab:intro-rq-mapping}。

\begin{table}[htbp]
\centering
\caption{研究问题、代表性组合与章节的对应关系}
\label{tab:intro-rq-mapping}
\thesistablefont
\setlength{\tabcolsep}{3pt}
\renewcommand{\arraystretch}{1.18}
\begin{tabularx}{\textwidth}{@{}p{.8cm}p{3.0cm}Xp{2.5cm}@{}}
\toprule
问题 & 代表性组合 & 研究内容与实验范围 & 对应位置 \\
\midrule
RQ1 & 多任务组织与平台设计 & 设计任务与阶段配置、模型或预测提交、评分与可视化流程；形成接口与测试方案 & 第\ref{chap:benchmark}章 \\
RQ2 & 多维评价与分割基准 & 定义共同评价指标；在域增量、类别增量和跨器官分割基准中分析学习、保持、泛化与资源条件 & 第\ref{chap:medcl}章指标；第\ref{chap:segmentation}章基准 \\
RQ3 & 分割／域与类别增量／弱监督／有限回放 & 沿用分割基准的六中心和三任务类别序列评价 ZSDERpp 及其背景适配 & 第\ref{chap:scribblecl}章 \\
RQ4 & 分类／类别扩展／联邦／无历史原始样本回放 & 研究跨客户端子空间融合；采用任务身份已知的多头分类评价 & 第\ref{chap:fedsubmerge}章 \\
RQ5 & 配准／任务增量／有限经验回放 & 以 SAMCL 在固定跨器官与跨模态任务序列中研究历史保持、当前适应和前向泛化 & 第\ref{chap:registration}章 \\
\bottomrule
\end{tabularx}
\end{table}

图~\ref{fig:thesis-work-logic}展示全文结构。MedCL 提供任务、数据和阶段结果的组织方式，第四至第六章据此开展分割、分类和配准实验。它们分别关注稀疏监督、无历史原始样本回放和有限图像对回放，采用 Dice、准确率或 TRE 评价任务性能，再比较新任务学习、旧任务保持及所需资源。

\begin{figure}[htbp]
\centering
\resizebox{\thesisfigurewidth}{!}{% Top-to-bottom thesis structure with parallel research branches in grayscale.
\begingroup
\linespread{1}\selectfont
\begin{tikzpicture}[
  font=\thesisfigurefont,
  box/.style={draw=black!55,rounded corners=2pt,align=center,inner sep=4pt,line width=.6pt},
  wide/.style={box,minimum width=13.1cm,fill=black!5,minimum height=.8cm},
  branch/.style={box,minimum width=3.9cm,minimum height=2.7cm,fill=white},
  flow/.style={-{Latex[length=2mm]},draw=black!65,line width=.65pt},
  node distance=5mm
]
\node[wide] (problem) {\textbf{共同问题：任务持续演化与训练信息受限}};
\node[wide,below=of problem,minimum height=1.05cm] (framework)
  {\textbf{第三章：MedCL 评测平台设计与分割基准}\\数据与任务设定\quad 阶段矩阵\quad 多维评价指标};
\node[branch,below=1.05cm of framework] (fed)
  {\textbf{第五章：类别扩展}\\\textbf{联邦持续分类}\\[5pt]无原始样本回放\\子空间融合与梯度保护};
\node[branch,left=5mm of fed] (seg)
  {\textbf{第四章：域与类别增量}\\\textbf{弱监督持续分割}\\[5pt]有限回放\enspace ·\enspace 涂鸦监督\\监督增强与特征一致性};
\node[branch,right=5mm of fed] (reg)
  {\textbf{第六章：任务增量}\\\textbf{持续医学图像配准}\\[5pt]有限图像对回放\\元学习与锐度感知优化};
\node[wide,below=6mm of fed,minimum height=1.05cm,fill=black!8] (evaluation)
  {\textbf{共同评价：任务性能与持续学习能力}\\阶段矩阵\quad BWTR\quad RMA\quad MPE\quad DRR};
\node[wide,below=of evaluation] (future)
  {\textbf{第七章：总结与展望}\qquad 多模态持续学习};
\draw[flow] (problem) -- (framework);
\draw[flow] (framework.south -| seg.north) -- node[right,align=left] {数据与评价\\直接承接} (seg.north);
\draw[flow] (framework.south) -- node[right] {阶段矩阵与指标} (fed.north);
\draw[flow] (framework.south -| reg.north) -- (reg.north);
\draw[flow] (seg.south) -- (evaluation.north -| seg.south);
\draw[flow] (fed.south) -- (evaluation.north);
\draw[flow] (reg.south) -- (evaluation.north -| reg.south);
\draw[flow] (evaluation) -- (future);
\end{tikzpicture}
\endgroup
}
\caption[本文研究工作的总体技术路线]{本文研究工作的总体技术路线。第三章介绍 MedCL 评测平台与多维指标；第四章构建持续分割基准并研究弱监督方法，第五、六章分别研究类别扩展下的联邦持续分类和跨器官、跨模态持续配准。三章使用共同的场景定义与阶段矩阵，并分别评价历史保持、当前学习及资源使用。}
\label{fig:thesis-work-logic}
\end{figure}

\textbf{（1）MedCL 医学影像持续学习评测平台设计。}
第三章设计面向分割、分类和配准的 MedCL 评测平台，明确任务顺序、阶段模型或预测、测试设置及多维评分之间的对应关系。用户提交模型或预测，平台负责检查、调用、评分和结果展示，不承担持续学习训练。平台以阶段--任务矩阵联系任务性能与持续学习能力，定义历史保持、当前学习、前向泛化及模型和数据资源指标，为后续三类任务提供共同的评价基础。

\textbf{（2）持续医学图像分割基准与弱监督方法。}
第四章首先构建域增量、类别增量和跨器官任务增量分割基准，在统一病例划分、任务顺序、网络与训练预算下比较代表性方法，分析学习、保持与资源之间的取舍\cite{wang2026benchmark}。随后，固定六中心前列腺 Domain-CL 和三阶段心脏结构 Class-CL 的数据与测试设置，将训练标注改为模拟涂鸦。ZScribbleSeg 提供静态弱监督基础，ScribbleCL 表示持续学习研究问题，ZSDERpp 结合全局一致性、空间先验、图像--涂鸦回放和参考特征约束处理监督不足与遗忘\cite{zhang2026zscribbleseg}，并在类别增量中结合 MiB 处理背景语义变化。密集监督与弱监督实验共同使用平台指标，分别以 E-FWT 和 WCD 补充域增量与类别增量评价。

\textbf{（3）类别扩展下的无回放联邦医学影像持续分类。}
第五章研究医学分类类别持续扩展、原始影像位于多个逻辑客户端且历史原始样本不回放时的协同学习。FedSubMerge 在客户端递归构建低秩主梯度子空间，服务器进行跨客户端融合，后续本地训练将梯度投影到保护空间的正交补；FedSubMerge-AD 根据逐层子空间距离选择相关客户端。实验采用已知测试任务身份的任务特定分类头，以准确率构建 MedCL 阶段矩阵，结合 ACC、BWTR、RMA、MPE、DRR 和通信开销分析方法表现。

\textbf{（4）任务增量医学影像持续配准。}
第六章研究跨器官、跨模态和不同配准关系顺序到达时的持续更新。SAMCL 将容量受限的历史图像对回放、一阶元持续学习和锐度感知优化结合，在四任务序列中评价配准精度、历史保持和当前适应\cite{wang2024samcl}。实验沿用 MedCL 的阶段矩阵，分别报告 Dice 和 TRE 的最终性能、后向变化、BWTR 与 RMA，并通过 MPE、DRR 和回放容量分析资源使用。

\section{本文主要创新与贡献}
\label{sec:intro-contributions}

围绕上述研究内容，本文的主要创新与贡献概括如下。

\textbf{（1）设计面向多类医学影像分析任务的 MedCL 持续学习评测平台。}
本文区分任务类型、增量设定和训练条件，设计阶段模型或预测的提交、评分与展示流程。平台通过阶段--任务矩阵比较新旧任务表现，并为联邦结果预留客户端分组评价。第三章给出需求、架构、接口、交互和测试方案，其阶段记录及历史保持、当前学习和资源指标供后续分割、分类与配准实验使用。

\textbf{（2）建立持续分割基准，并提出监督增强与特征一致性约束相结合的弱监督方法。}
本文构建域增量、类别增量和跨器官任务增量分割基准，在统一实验条件下比较代表性方法，表明较低遗忘并不必然带来更好的新任务学习或未来域泛化\cite{wang2026benchmark}。在 ZScribbleSeg 的监督增强基础上，进一步结合有限图像--涂鸦回放与参考特征约束，提出 ZSDERpp。六中心域增量实验中，其 A-Dice 为 $0.7687$、BWTR 为 $-0.1520$；三任务类别增量实验中，与 MiB 结合后的 A-Dice 为 $0.7793$、WCD 为 $0.7292$。两类实验沿用基准的数据划分和测试参照，分析稀疏监督下的学习与保持能力。

\textbf{（3）提出基于跨客户端主梯度子空间融合的无回放联邦持续分类方法。}
FedSubMerge 通过奇异值加权融合汇集跨客户端历史方向，并以正交补投影约束后续更新；FedSubMerge-AD 按层选择相关客户端，以减轻不相关方向对当前学习的限制。本文分析局部损失变化、两阶段子空间截断及通信存储开销，并在七个联邦设置中比较两种融合方式的分类性能与历史保持。

\textbf{（4）提出面向跨器官与跨模态任务变化的锐度感知元经验回放持续配准方法。}
SAMCL 将有限经验回放、一阶元持续学习与锐度感知更新结合，用于跨器官和跨模态的连续形变估计。四任务实验比较最终配准精度、历史保持和回放容量的影响，并在预设测试方向上分析模型的直接泛化表现。

\section{论文组织结构}
\label{sec:intro-organization}

全文共七章。第一章介绍研究背景、相关工作和主要问题。第二章梳理分割、分类与配准基础，以及持续学习的设定和关键技术。第三章设计 MedCL 评测平台，说明任务与数据组织、阶段矩阵和评价指标。第四章先进行密集监督持续分割基准评测，再在相同数据划分与测试设置下研究弱监督分割。第五章研究无回放联邦持续分类，第六章研究有限回放持续配准。第四至第六章均沿用第三章的阶段评价方式。第七章总结全文工作与局限，并讨论持续预训练、参数高效持续适配和持续知识编辑等后续方向。